\chapter{Semantics, Pragmatics, and Usage}\label{ch:semantics}

% ============================================================
% STATUS: SCAFFOLD — port and revise from
% archive/2020/chapters/08-semantics-and-usage.tex.
% ============================================================

\section{Register and formality}\label{sec:register}

% [Port from 2020 archive; update case terminology in examples]

\section{Politeness and the T-V distinction}\label{sec:politeness}

% ============================================================
% The 2020 archive had a detailed treatment of the T-V distinction:
%   - 2pl used as polite singular
%   - 1pl used for speaker humility in formal contexts
%   - Pro-drop as a politeness strategy (avoiding pronouns entirely)
%   - Use of names and titles as quasi-honorifics
% The plural marker ne/ně as an honorific with proper names and
% kinship terms (ne Stám Kovárž, ne mlazka) should be discussed here.
% ============================================================

\section{Honorifics}\label{sec:honorifics}

% [Develop from 2020 archive notes on ně/ne as honorific particle]

\section{Phatic expressions}\label{sec:phatic}

% [Greetings, leave-takings, thanks, apologies — cross-reference with
%  §\,\ref{sec:fixed-converbs} in Chapter~\ref{ch:complex}]

\section{Idiomatic expressions}\label{sec:idioms}

% [Selected idioms from the 2020 archive, updated glosses]

